%======================================================================
\section{Introduction}\label{sec:intro}
%======================================================================

In a two-player nonlocal game a referee sends a question to each of two players, who cannot communicate, and they win if their answers
satisfy a fixed predicate. Players who share an entangled state can sometimes win with a probability that no classical strategy reaches,
and for some games the best winning probability, the \emph{quantum value} $\omq(G)$, is approached only by strategies of ever larger
dimension. For such a game the basic quantitative question is how much entanglement near-optimal play needs. Lower bounds on it are
dimension witnesses: from the winning statistics alone, they certify that the players share a large entangled state. Write $\Ent(G,\eps)$
for the least local dimension of a strategy that wins with probability at least $\omq(G)-\eps$. How fast does $\Ent(G,\eps)$ grow as $\eps\to0$?

The question has a short history. For two-player games of quantum value one, the first explicit examples gave polynomial growth.
Slofstra and Vidick built a linear system game with $c\,\eps^{-1/6}\le\Ent(G,\eps)\le C\eps^{-1/2}$ and wrote that they did not expect the
inverse polynomial scaling to be optimal~\cite{SV18}. Slofstra then reached $\exp(\eps^{-c})$ for a game of commuting-operator value
one, and noted that it is not clear whether its quantum value is one, in part because his group is not known to be hyperlinear and in
part because of limitations of the embedding theorem of~\cite{Slo20} that he used~\cite{Slo18}. Games with super-polynomial growth
are known with three players~\cite{JLV20}, and for two players when the value is not one~\cite{Col20}. There is also a uniform obstruction.
Deciding whether $\omq(G)=1$ is $\Pi_2$-complete~\cite[Thm.~1.1]{MNY22}, while for fixed $D$ and $k$ it is decidable, by
Tarski--Seidenberg, whether some strategy of local dimension at most $D$ wins with probability at least $1-1/k$; so no total computable
function of the game and the precision bounds $\Ent(G,\eps)$ for all games of quantum value one. This names no game and gives no rate
for any one game. Dykema, Paulsen and Prakash showed, through correlation functions of graphs, that
the set of quantum correlations is not closed already for five questions and two answers per player~\cite{DPP19}, so that approximating
some limits of finite-dimensional correlations needs unbounded dimension. \Cref{tab:known}
collects the explicit bounds.

\subsection{The result}

We give an explicit two-player game of quantum value one in which every near-optimal strategy needs exponential entanglement, and we
determine the exponent. Throughout, $\bfun(\eps)=2^{-4(24+\log\log(1/\eps))^2}$; this is $\eps^{o(1)}$ as $\eps\to0$. Logarithms are to
base two.

\begin{theoremA}
There is a binary linear system game $G$, with $1687$ equations in $1993$ variables and at most three variables in each equation, such that
$\omq(G)=1$ and no finite-dimensional strategy wins $G$ with certainty, and there are explicit constants $c,C,\eps_0>0$ with the
following properties for $0<\eps\le\eps_0$.
\begin{enumerate}[label=(\roman*),itemsep=1pt]
\item Every strategy that wins $G$ with probability at least $1-\eps$ has entanglement entropy
\[
S(\rho_B)\ \ge\ c\,\eps^{-1}\,\bfun(\eps),
\]
and hence Schmidt rank, and local dimension, at least $2^{c\,\eps^{-1}\bfun(\eps)}$. For $0<\eps\le\eps_0/2$ the same lower bound,
with $c/4$ in place of $c$, holds for the entanglement of formation of every mixed strategy.
\item Some strategy that wins with probability at least $1-\eps$ uses a maximally entangled state of local dimension at most
$\exp\bigl(C\eps^{-1}\log(1/\eps)\bigr)$.
\end{enumerate}
In particular $\log\Ent(G,\eps)=\eps^{-1+o(1)}$, and the same holds for the least entanglement entropy and for strategies on maximally
entangled states.
\end{theoremA}

The precise statement is \cref{thm:main}, with $c=4.9\cdot10^{-34}$, $C=7.4\cdot10^{10}$ and $\eps_0=2^{-5500}$; these values are far
from optimal. \Cref{box:game} describes the game. To our knowledge $G$ is the first explicit two-player game
of quantum value one that needs more than polynomial entanglement, and the first two-player game of value one whose entanglement is pinned
to a scale $\exp(\eps^{-\Theta(1)})$. The lower bound is on the entanglement \emph{entropy}, so it counts ebits and not only dimensions:
near-optimal play of $G$ consumes $\eps^{-1+o(1)}$ ebits, whatever the shape of the shared state.

\begin{table}[t]
\centering
\footnotesize\setlength{\tabcolsep}{3.5pt}%
\begin{tabular}{@{}llllll@{}}
\toprule
 & players & value & needed & sufficient & remark \\
\midrule
\cite{SV18} & 2 & $\omq=1$ & $c\,\eps^{-1/6}$ & $C\eps^{-1/2}$ & linear system game \\
\cite{MV15} & 2 & $1$ & $c\,\eps^{-1/2}$ & $\eps^{-O(1)}$ & answers of unbounded length \\
\cite{Slo18} & 2 & $\omqc=1$ & $\exp(\eps^{-c})$ & & $\omq=1$ not known \\
\cite{Col20} & 2 & below $1$ & $2^{\Omega(\eps^{-1/8})}$ & $2^{O(\eps^{-1})}$ & within $\eps$ of the value \\
\cite{JLV20} & 3 & $1$ & $2^{\Omega(\eps^{-c})}$ & $2^{O(\eps^{-1})}$ & three players \\
\midrule
Theorem~A & 2 & $\omq=1$ & $\exp(\eps^{-1+o(1)})$ & $\exp(C\eps^{-1}\log\frac1\eps)$ & linear system game; also in ebits \\
\bottomrule
\end{tabular}
\caption{Local dimension needed, and known to suffice, to win with probability $\omega-\eps$, for games whose value $\omega$ is not attained
by finite-dimensional strategies. For Coladangelo's game the bounds concern coming within $\eps$ of its value.
Gao, Yu and Zhi recently gave an explicit linear system game whose commuting-operator
value is one and whose quantum value is below one~\cite{GYZ26}, a different question.}
\label{tab:known}
\end{table}

\paragraph{The group behind the game.} A linear system game is governed by its solution group $\Gamma$, which contains a central involution
$J$: a near-optimal strategy is, up to small errors, a representation of $\Gamma$ in which $J$ acts as $-I$, and the entanglement of the
game is the price of keeping $J$ away from the identity in approximate representations~\cite{SV18}. So we need a group in which $J$ survives
in approximate representations, which gives value one, but only in very large dimension.

Our group $\GP$ is a lamplighter whose lamps sit on configurations rather than on points (\cref{sec:group}, \cref{fig:lamps}). The points
are those of the three rays of Houghton's group $H_3$; a configuration is a finite set of points, written as a binary vector; and the
lighter is a group of affine maps of configurations that permute points, toggle a point, or add one point to another. Each lamp is a Pauli
pair $x_v,z_v$, a copy of the dihedral group of order eight, and the centres $(x_vz_v)^2$ of all lamps are one central involution $J$. On
$N$ points there are $2^N$ configurations, hence $2^N$ Pauli pairs. Yet the generators of the lighter either move points along the rays,
as Houghton's group does, or involve at most two points of a configuration, so any two lamps can be brought together cheaply: their
commutator has area $2^{O(\log^2N)}$ (\cref{thm:pair-area}),
priced by a quasi-polynomial Dehn bound for $H_3$ that we prove in \cref{app:dehn}. The group is elementary amenable, hence hyperlinear,
which removes the first of Slofstra's obstacles. For the second we give explicitly a presentation of the type that the Slofstra--Vidick
reduction needs, with one non-involutive generator, whose $29$ derivations are checked by machine (\cref{sec:game}).

% Figure 1: the three rays, a configuration, and a transporter (introduction).
\begin{figure}[t]
\centering
\begin{tikzpicture}[x=1cm,y=1cm,font=\small,>=Stealth,baseline=(current bounding box.center),
  pt/.style={circle,draw,fill=white,inner sep=0pt,minimum size=5pt},
  on/.style={pt,fill=black},
  lbl/.style={font=\footnotesize,inner sep=1pt}]
\def\u{0.7}
\foreach \i/\ang in {1/0,2/120,3/240}{
  \draw[black!50] (0,0) -- (\ang:4.4*\u);
  \node at (\ang:4.85*\u) {$\cdots$};
  \foreach \j in {1,...,4}{\node[pt] (p\i\j) at (\ang:\j*\u) {};}
}
\node[lbl,anchor=west] at ($(120:4.8*\u)+(0.3,0)$) {ray $2$};
\node[lbl,anchor=west] at ($(240:4.8*\u)+(0.3,0)$) {ray $3$};
\node[lbl,anchor=north] at ($(0:4.8*\u)+(0.1,-0.2)$) {ray $1$};
\foreach \j in {1,...,4}{\node[lbl,anchor=north] at ($(p1\j)+(0,-0.16)$) {$\j$};}
\node[on] at (p11) {};
\node[on] at (p13) {};
\node[lbl,anchor=south] at ($(p12)+(0,0.3)$) {$u=e_1+e_3$};
\end{tikzpicture}%
\hspace{1.2cm}%
\begin{tikzpicture}[x=0.85cm,y=0.85cm,font=\small,>=Stealth,baseline=(current bounding box.center),
  lamp/.style={circle,draw,fill=white,inner sep=0pt,minimum size=7pt},
  mv/.style={blue!50!black,very thick,->,shorten >=4pt,shorten <=4pt},
  lbl/.style={font=\footnotesize,inner sep=1pt}]
\coordinate (v0) at (0,0);
\coordinate (v1) at (2.4,0);
\coordinate (v2) at (0,2.4);
\coordinate (v12) at (2.4,2.4);
\coordinate (v3) at (1.3,1.1);
\coordinate (v13) at (3.7,1.1);
\coordinate (v23) at (1.3,3.5);
\coordinate (v123) at (3.7,3.5);
\draw[black!40] (v0) -- (v1) -- (v12) -- (v2) -- cycle;
\draw[black!40] (v3) -- (v13) -- (v123) -- (v23) -- cycle;
\draw[black!40] (v0) -- (v3) (v1) -- (v13) (v2) -- (v23) (v12) -- (v123);
\draw[mv] (v0) -- (v3);
\draw[mv] (v3) -- (v13);
\node[lbl,blue!50!black,anchor=south east] at ($(v0)!0.5!(v3)+(-0.02,0.06)$) {$\tau_3$};
\node[lbl,blue!50!black,anchor=south] at ($(v3)!0.3!(v13)+(0,0.08)$) {$\tau_1$};
\foreach \v in {v0,v1,v2,v12,v3,v13,v23,v123}{\node[lamp] at (\v) {};}
\node[lbl,anchor=north east] at ($(v0)+(-0.1,-0.1)$) {$0$: $x,z$};
\node[lbl,anchor=north west] at ($(v13)+(0.1,-0.12)$) {$u$: $x_u,z_u$};
\node[lbl,anchor=north] at ($(v1)+(0,-0.18)$) {$e_1$};
\node[lbl,anchor=east] at ($(v2)+(-0.18,0)$) {$e_2$};
\node[lbl,anchor=south east] at ($(v3)+(-0.1,0.08)$) {$e_3$};
\end{tikzpicture}
\caption{Left: the points of the three rays of Houghton's group and a configuration $u=e_1+e_3$, a finite set of points (filled). Right: the
$2^3$ configurations supported on the points $1,2,3$ of ray~$1$, each carrying a Pauli pair (a circle), with $(x_vz_v)^2=J$ for all $v$. The
transporter $T_u=\tau_1\tau_3$, a word in the generators, carries the configuration $0$ to $u$ (bold path), and conjugation by it carries the
pair $x,z$ at $0$ to the pair $x_u=T_uxT_u^{-1}$, $z_u=T_uzT_u^{-1}$ at $u$. On $N$ points there are $2^N$ pairs, and on the sector where $J$
acts as $-1$ each acts as a qubit.}
\label{fig:lamps}
\end{figure}


\subsection{Why the state must be kept}

A strategy that wins with probability $1-\eps$ controls the players' operators to precision about $\sqrt\eps$, measured in their shared
state: if $\rho$ is Bob's reduced density matrix and $\xi=\rho^{1/2}$, his observables satisfy the relations of the solution group up to
$\sqrt\eps$ in the norm $\hs{Y\xi}=\Tr(Y^*Y\rho)^{1/2}$, and they nearly commute with $\xi$, again to order $\sqrt\eps$. The usual route to a
lower bound then \emph{flattens} the state: it passes to a spectral projection of $\rho$~\cite{SV18}, to a strategy on a maximally entangled
state~\cite{Pad22}, or to the operator norm. Each of these steps costs a power of $\eps$. For linear system games the precision $\sqrt\eps$
becomes $\eps^{1/4}$~\cite[Cor.~5.2]{SV18}; for our game this would give only $\log\Ent(G,\eps)\ge\eps^{-1/2+o(1)}$, the exponent of
maximally entangled strategies halved.

We avoid the loss by never leaving the state. The counting argument below works with the Hilbert--Schmidt vector $\xi$ itself: it
compares quantities measured in one and the same state, and it pays for every change of state with a \emph{state commutator}
$\hs{[Y,\xi]}$, which near-optimal play supplies at precision $\sqrt\eps$. Two results make this precise, and both may be of independent use.

The first is a robust form of the fact that $M$ pairs of anticommuting involutions need $M$ qubits, valid in an arbitrary state, with the
errors of different pairs adding in squares rather than in norms.

\begin{theoremC}[\cref{thm:count,cor:count}]
Let $X_1,\dots,X_M,Z_1,\dots,Z_M$ be involutions on $\C^d$, $M\ge2$, let $\rho$ be a density matrix and $\xi=\rho^{1/2}$. If for all $i\ne j$
\[
\hs{[Z_i,Z_j]\xi},\ \ \hs{[X_i,Z_j]\xi},\ \ \hs{\{X_i,Z_i\}\xi},\ \ \hs{[X_i,\xi]},\ \ \hs{[Z_i,\xi]}\ \le\ \kappa
\]
and $16M\kappa^2\le\frac1{10}$, then $S(\rho)\ge\frac34M$.
\end{theoremC}

The full statement, \cref{thm:count}, bounds $S(\rho)$ from below by a sum over the pairs in which the errors of each pair enter only through sums of
squares. The state commutators cannot be dropped (\cref{ex:commutators-needed}). Stability theorems of Gowers--Hatami type~\cite{GH17} and the
many-qubit tests of Natarajan and Vidick~\cite{NV17} need control of a whole Pauli group; \cref{thm:count} needs only the pairwise relations,
and the $\ell^2$ tolerance is what produces the exponent $1$ in Theorem~A. Pairwise approximate relations with a fixed error per pair allow many
approximate qubits to overlap in a small space~\cite{CRSV17}; \cref{thm:count} is consistent with this, because it charges the errors of all
pairs together, through a sum of squares.

The second result says that for linear system games nothing is lost in either direction when the state is kept. For a finite presentation
and a word $w$, let $\hlps(w;e)$ be the least dimension of a homomorphism $\phi$ from the free group to $\U(d)$ together with a density
matrix $\rho$ such that every relator is $e$-close to the identity in the norm $\hs{(\cdot)\xi}$, every generator has state commutator at
most $e$, and $\hs{(\phi(w)+I)\xi}\le e$ (\cref{def:state-profile}).

\begin{theoremD}[\cref{thm:transfer,thm:return,cor:two-sided}]
Let $G$ be a binary linear system game of quantum value one with solution group $\Gamma$ and central element $J$. There are constants
$C,C'>0$, depending only on the number of equations, the number of variables per equation and the question distribution, such that for
every $\eps>0$
\[
\hlps_\Gamma\bigl(J;C\sqrt\eps\bigr)\ \le\ \Ent(G,\eps)\ \le\ \hlps_\Gamma\bigl(J;C'\sqrt\eps\bigr).
\]
Both directions keep the state: the lower bound uses Bob's reduced state, and the upper bound is realised on the canonical purification of
the state of the model. The same holds for the least entanglement entropy.
\end{theoremD}

For flat models the right-hand inequality is the classical direction of the correspondence of Slofstra and
Vidick~\cite[Prop.~4.4(a)]{SV18}. The left-hand one is their passage from strategies to approximate representations~\cite[Prop.~5.4]{SV18},
stopped before the state is flattened.

\subsection{The profile of the central involution}

On the side of the group, the lower bound of Theorem~A is a statement about approximate representations of $\GP$ in which $J$ is visible.
Let $\hlp(J;\delta,\eps)$ be the least dimension of an $\eps$-representation $\phi$ with $\fn{\phi(J)-I}\ge\delta$, the \emph{hyperlinear
profile} of $J$~\cite{SV18}, and let $\RP$ be the finite certificate of $27$ relations of $\GP$ of \cref{def:RP}.

\begin{theoremB}[\cref{thm:profile}]
Let $R\supseteq\RP$ be a finite set of words in the generators $p,q,t,d,x,z$ of $\GP$ that represent $1$ in $\GP$, and $J=(xz)^2$. There are
constants $c,C>0$, with $C$ depending on the lengths of the words of $R$, such that for $0<\delta\le2$ and $0<\eps\le2^{-1071}\delta$
\[
c\,\delta^2\eps^{-2}\bfun(\eps)\ \le\ \log\hlp_R(J;\delta,\eps)\ \le\ C\,\delta^2\eps^{-2}\log(1/\eps).
\]
The same two-sided bound $\log\hlps_R(J;e)=e^{-2+o(1)}$ holds for the state-weighted profile.
\end{theoremB}

So $\log\hlp(J;\delta,\eps)=\eps^{-2+o(1)}$. By~\cite[Prop.~7.1]{DM-lamps} the group $\GP$ is finitely presented (nothing here uses
this; see \cref{rem:GP-fp}), so it is, to our knowledge, the first finitely presented amenable group with a central involution of
superpolynomial hyperlinear profile. The upper bound comes from \emph{hash models} (\cref{sec:upper}): we keep only the
lamps that a short word lights, on a register of qubits indexed by a random affine hash of the configurations, and a clock handles the
abelian quotient of the lighter. Slofstra's group in~\cite{Slo18} has $\log\hlp\ge c\,\eps^{-a}$ for every $a<\frac12$, and is not known to be
hyperlinear; the sofic group of Slofstra and Vidick has a central involution of polynomial profile~\cite[Prop.~3.2]{SV18}.

\subsection{Two sharp rates}

The same two tools close the gaps in two classical examples, in which the known lower bounds lost a power of $\eps$ by flattening.

The XOR games CHSH$(n)$ of Slofstra~\cite{Slo11} force $n$ pairwise anticommuting observables. Ostrev showed that every nearly optimal
strategy approximately contains the algebraic structure of the canonical optimal one~\cite{Ost16}; the question here is how much
entanglement that forces as a function of the error. Ostrev and Vidick showed that a strategy within
a factor $1-\eps$ of the optimal bias needs $\Omega(\eps^{-1/5})$ ebits, using $n=\Theta(\eps^{-1/5})$, and that every XOR game has an
$\eps$-optimal strategy with $\lceil1/\eps\rceil$ ebits~\cite{OV18}.

\begin{theoremE}[\cref{thm:chsh,cor:chsh-ebits,cor:chsh-ef}]
For $n\ge3$ and $0<\eps\le1$, the least entanglement entropy of a strategy for CHSH$(n)$ whose bias is at least $1-\eps$ times the optimal
bias lies between $10^{-7}\min(n,1/\eps)/6-\frac12$ and $\min(\lfloor n/2\rfloor,\lceil1/\eps\rceil)$; so it is $\Theta(\min(n,1/\eps))$ as
$\min(n,1/\eps)\to\infty$. The least entanglement of formation of a mixed strategy obeys the same bounds, with $10^{-7}/24$ and $\frac14$
in place of $10^{-7}/6$ and $\frac12$.
\end{theoremE}

This improves the exponent $\frac15$ of~\cite{OV18} to the optimal $1$. Slofstra and Vidick built a linear system game $\Gsv$ from a group $K$
whose central involution $c$ has hyperlinear profile between $\eps^{-2/3}$ and $\eps^{-1}$, and gave the bounds
$c\,\eps^{-1/6}\le\Ent(\Gsv,\eps)\le C\eps^{-1/2}$~\cite{SV18}.

\begin{theoremF}[\cref{cor:K-profile,thm:sv-game}]
The hyperlinear profile of $c$ in $K$ is $\Theta(1/\eps)$, for every fixed $\delta$, and $\Ent(\Gsv,\eps)=\Theta(\eps^{-1/2})$.
\end{theoremF}

Here the relations make the spectral measure of one generator, measured in the state, nearly invariant under doubling and under a half-turn
of the circle; averaging over the two inverse branches of the doubling map contracts the Wasserstein-$1$ distance by a factor two, so the
measure is within $\frac{29\pi}2$ times the precision of Haar measure (\cref{prop:K-W1}), while a measure with $d$ atoms is at distance at
least $\pi/(2d)$.

The upper bound $\Ent(\Gsv,\eps)\le C\eps^{-1/2}$ is that of Slofstra and Vidick~\cite[Thm.~1.1]{SV18}. Its proof, through that
of~\cite[Prop.~6.2]{SV18}, lifts approximate representations of $K$ to the solution group by~\cite[Prop.~4.8]{Slo19}, which removes one
non-involutive generator at a time; $K$ has two, and after both removals the image of $c$ is no longer balanced, so the last step of that
proof does not apply as written (\cref{sec:rates-sv}). Only the upper bound is affected: the value one of the game, the lower bounds and
the nontriviality of $J$ do not use this step, and with at most one non-involutive generator the argument is complete. For the explicit
models of $K$ we give a lift that keeps $c$ balanced (\cref{lem:sv-lift}), so~\cite[Thm.~1.1]{SV18} holds as stated.

\subsection{How the proof of Theorem A works}

We describe the lower bound; the upper bound runs the chain backwards using flat approximate models in which the marked involution is
exactly $-I$. \Cref{game:fig-map} in \cref{sec:game}
shows both chains with their constants.

\emph{From the game to a state-weighted model.} A strategy that wins with probability $1-\eps$ gives, on the support of Bob's reduced state
$\rho$, involutions for the variables of the game that satisfy the relations of the solution group $\Gamma$, and commute with $\xi$, up to
$O(\sqrt\eps)$ in the state (\cref{thm:transfer}). This is bookkeeping, but it is the step at which earlier arguments flattened the state.

\emph{From the game to the group.} The game is built from a presentation $\Pres$ of Slofstra's type $E\Gamma_0$ through Slofstra's
embedding, run so that the distinguished involution goes to $-I$ (\cref{app:gadget}), and the $27$ relations of the certificate $\RP$ are
consequences of $\Pres$ of area at most $60$ (\cref{prop:presentation}). So the model of $\Gamma$ restricts to an assignment of unitaries to
the generators $p,q,t,d,x,z$ of $\GP$ satisfying $\RP$ to $O(\sqrt\eps)$ in the state, in which the image of $J$ stays at distance about one
from the identity.

\emph{Lamps on the sector.} On the sector where (the rounded image of) $J$ acts as $-1$, the transported lamps $x_u,z_u$, $u\subseteq\{1,\dots,N\}$,
are $2^N$ approximate Pauli pairs measured in the state: within a pair they nearly anticommute because $(x_uz_u)^2=J=-1$, and across pairs
they nearly commute at a cost of the pair area $\Lambda_N=2^{O(\log^2N)}$ times the precision (\cref{thm:pair-area}). Charging conjugators
through state commutators squares this cost.

\emph{Counting.} Theorem~C applies with $M=2^N$ and $\kappa\approx\Lambda_N^2\sqrt\eps$, as long as $2^N\Lambda_N^4\eps\lesssim1$. Choosing the
largest such $N$ gives $2^N=\eps^{-1+o(1)}$ qubits of entropy on the sector (\cref{thm:sector}), and the sector carries a constant fraction of
the state (\cref{lem:sector-entropy}). This is where the exponent $1$ comes from: precision $\sqrt\eps$, errors adding in squares.

\emph{The upper bound.} The hash models give $\eps$-representations of $\GPp$ with $J\mapsto-I$ exactly in dimension
$\exp(C\eps^{-2}\log(1/\eps))$; they give approximate representations of $\Pres$ with $w\mapsto-I$ and then of $\Gamma$ with $J\mapsto-I$
(the embedding theorem, \cref{thm:embedding} in \cref{app:gadget}), and the return direction of
Theorem~D turns them into strategies on a maximally entangled state that lose with probability $O(\eps^2)$ (\cref{thm:return}).

\subsection{Related work and what is new}

The architecture of the group follows Slofstra's~\cite{Slo18}: a central involution $J$, a family of involutions whose commutators are
$J$ or $1$, exponentially many of them reached by short words, and a transfer of approximate representations along the group.
In~\cite{Slo18} the involutions generate, with $J$, an infinite Clifford group, any two of them anticommuting up to $J$, and they are
reached through the HNN doubling $tzt^{-1}=z^2$; here they form Pauli pairs and are reached through Houghton's group acting on
configurations, which keeps the group amenable. The passage from groups to games is Slofstra's embedding~\cite{Slo19,Slo20} in the explicit form of Slofstra and
Vidick~\cite{SV18}, and the identification of the entanglement of a linear system game with the hyperlinear profile of $J$ is theirs. The
pair-area argument and the hash models adapt our earlier construction for lamps that are copies of $S_3$, where they bound the sofic profile
of an amenable group~\cite{DM-lamps}; nothing in this paper depends on that preprint.

On the analytic side, the Hilbert--Schmidt picture of a shared state, with the reduced state as a factor $\xi$, is the standard form of the
state. Slofstra and Vidick use it for general strategies~\cite{SV18}, and Marrakchi and de la Salle use it to round near-synchronous
strategies~\cite{MdlS23}. The coupling of two spectral measures in one state is Connes' joint distribution of left and right
multiplication~\cite{Con76}. Robust tests of many qubits in the players' state are due to Natarajan and Vidick~\cite{NV17}, through a test
that samples the whole Pauli group; Pauli stability for synchronous strategies is in~\cite{dlS22}; state-dependent anticommutators appear in
the entropic uncertainty relations of Kaniewski, Tomamichel and Wehner~\cite{KTW14}; and rounding to maximally entangled or tracial strategies
is studied in~\cite{Vid22,Pad22,Lin23,MdlS23}.

What is new is the following.
\begin{enumerate}[label=(\arabic*),itemsep=1pt]
\item An explicit two-player game of quantum value one whose near-optimal strategies need $\eps^{-1+o(1)}$ ebits, with a matching upper
bound (Theorem~A).
\item An elementary amenable group with a central involution whose hyperlinear profile is $\exp(\eps^{-2+o(1)})$, two-sided, for every
finite set of valid relations that contains $\RP$ (Theorem~B), and a self-contained proof, from Johnson's presentation, of a
quasi-polynomial Dehn bound for Houghton's group $H_3$ (\cref{app:dehn}); our preprint~\cite{DM-H3} proves a polynomial bound, which
this paper does not use.
\item A count of qubits from pairwise Pauli relations in an arbitrary state, with errors adding in squares (Theorem~C).
\item A correspondence between linear system games and state-weighted approximate representations that loses nothing in either direction
(Theorem~D).
\item The optimal exponents for CHSH$(n)$ and for the game and group of Slofstra and Vidick (Theorems~E and~F).
\end{enumerate}

\subsection{Organization and conventions}

\Cref{sec:prelim} fixes norms, approximate representations, state-weighted models and linear system games. \Cref{sec:state} develops a
calculus for words in unitaries measured in a fixed state, \cref{sec:count} proves Theorem~C, and \cref{sec:transfer} proves Theorem~D.
\Cref{sec:group} defines the group $\GP$, \cref{sec:pair} proves the pair areas, \cref{sec:sector} the lower bound of Theorem~B and the sector
theorem, and \cref{sec:upper} the hash models and Theorem~B. \Cref{sec:game} builds the game and proves Theorem~A, and \cref{sec:rates}
proves Theorems~E and~F. \Cref{sec:open} lists open questions. \Cref{app:dehn} proves the Dehn bound for $H_3$, \cref{app:gadget} gives
Slofstra's embedding with the distinguished involution sent to $-I$, and \cref{app:cert} describes the machine checks of the finite
certificates.

Logarithms are to base two unless written $\ln$; $h$ is the binary entropy and $S$ the von Neumann entropy in bits. For words and group
elements $[a,b]=aba^{-1}b^{-1}$; for operators $[S,T]=ST-TS$ and $\{S,T\}=ST+TS$.
